\section{Dasar Teori}
Pada bagian ini diberikan definisi, contoh, lema, dan teorema yang digunakan dalam penelitian.

\begin{definisi}[\cite{penulis2015buku}]\label{def:metrik}
  Diberikan himpunan tak kosong $X$. Fungsi $d:X\times X\to\R$ disebut metrik pada $X$ jika untuk setiap $x,y,z\in X$ berlaku
  \begin{enumerate}
    \item[(a)] $d(x,y)\geq 0$ dan $d(x,y)=0$ jika dan hanya jika $x=y$,
    \item[(b)] $d(x,y)=d(y,x)$,
    \item[(c)] $d(x,z)\leq d(x,y)+d(y,z)$.
  \end{enumerate}
\end{definisi}

\begin{contoh}\label{con:metrik}
  Fungsi $d(x,y)=|x-y|$ untuk setiap $x,y\in\R$ merupakan metrik pada $\R$.
\end{contoh}

\subsection{Contoh Subbab}
\begin{lema}\label{lem:contoh}
  Tuliskan pernyataan lema di sini.
\end{lema}
\begin{proof}
  Tuliskan bukti lema di sini.
\end{proof}

\begin{teorema}[\cite{contoh2024artikel}]\label{thm:contoh}
  Tuliskan pernyataan teorema di sini, misalnya
  \begin{equation}\label{eq:contoh}
    \norm{x+y}\leq\norm{x}+\norm{y}.
  \end{equation}
\end{teorema}

Berdasarkan \autoref{def:metrik}, \autoref{con:metrik}, dan \autoref{thm:contoh}, khususnya Persamaan \eqref{eq:contoh}, \dots
